\section{静电场中的电介质}
在本章的前两节我们已经研究了导体在静电场中的行为，在这一节中我们将研究所谓的电介质，也就是绝缘体在静电场中的效应，也许最初我们认为既然绝缘体不导电，在电场中就不应有任何效应，但是通过一个简单的实验就能发现事实并非如此，考虑下图。
\begin{Figure}{电介质对电容器的影响}
    \subfigure[未置入电介质的电容器]
    {
        \label{图:未置入电介质的电容器}
        \inputfigure[scale=0.65]{Chapter09C_Figure01}
    }\qquad\qquad\qquad
    \subfigure[置入电介质的电容器]
    {
        \label{图:置入电介质的电容器}
        \inputfigure[scale=0.65]{Chapter09C_Figure02}
    }
\end{Figure}
如图\ref{图:未置入电介质的电容器}所示，在电容器上充上一些电荷，此时测得的电压为$U_1$，现在我们在电容器的极板间塞入一块电介质，例如玻璃板或树脂片等，如图\ref{图:置入电介质的电容器}，此时测得的电压为$U_2$，而现在有这样的实验事实，表明这里$U_1>U_2$，即置入电介质后电容器两侧的电压减小了。

怎么会这样呢？在这个过程中，导体极板上的电荷量$\pm\sigma_f$并不变化，导体产生的场强$\vb*{E}_f$也不变，但是电容器的电压$U$发生了变换，由于匀强电场中电压$U=Ed$，因此电容器中的电场强度$\vb*{E}$必然发生了变化。所以唯一的可能就是，电介质参与了电场作用，电介质表面产生了一些电荷，靠近导体正极板处具有负电荷$-\sigma_b$，靠近导体负极板处具有正电荷$+\sigma_b$，这就在原有$\vb*{E}_f$的基础上产生了一个反向场强$\vb*{E}_b$，从而削弱了电场强度和电压，但由于电场强度没有被削弱至零，因此可以推断出电介质表面的电荷$\sigma_b$小于导体表面的电荷$\sigma_f$。

那么现在的问题就是，电介质为何会在电场中表现出这样的行为？

\subsection{电介质的极化}
在化学中，我们知道，分子整体是电中性的，但是由于原子间颠覆性存在差异，电子对会偏向或转移至电负性较强的原子，因此在分子中，某些原子带正电，某些原子带负电，这是一个很复杂的带电系统，但是由于我们考虑的电场的线度远大于分子自身的线度，如\Refx{图}{有极分子}和\Refx{图}{无极分子}所示，因此可以认为分子中的正负电荷是分别集中在某一点上的，这分别称为\uwave{正电荷的重心}和\uwave{负电荷的重心}，这两个重心未必是重合的。
\begin{Figure}{分子的电矩模型}
    \subfigure[有极分子]
    {
        \label{图:有极分子}
        \inputfigure[scale=0.9]{Chapter09C_Figure03}
    }\qquad\qquad
    \subfigure[无极分子]
    {
        \label{图:无极分子}
        \inputfigure[scale=0.9]{Chapter09C_Figure04}
    }
\end{Figure}
\begin{itemize}
    \item 如果分子的正负电荷的重心并不重合，那么就将这种分子称为\uwave{有极分子}，此时分子模型可以被简化为相距一定距离的的一对等量异号电荷，即电偶极子，具有一定的电偶极矩，称为\uwave{分子的固有电矩}，常见的有极分子包含\ce{H2O}、\ce{H2S}、\ce{SO2}、\ce{NH3}等。
    \item 如果分子的正负电荷的重心是重合的，那么就将这种分子称为\uwave{无极分子}，无极分子不具有固有电矩，常见的无极分子包含\ce{H2}、\ce{O2}、\ce{N2}、\ce{CO2}、\ce{CH4}、\ce{CCl4}等。
\end{itemize}
下面我们就来分别考虑，有极分子和无极分子构成的电介质是如何与电场发生相互作用的。

\subsubsection{无极分子的位移极化}
对于无极分子构成的均匀电介质，如\Refx{图}{无极分子的位移极化}所示。

当电介质外没有电场时，由于无极分子本身没有固有电矩，因此在宏观上不产生电场。

当电介质外存在电场时，由于电场力的作用，使得电子云的分布发生了一定的偏移，从而分子的正负电荷重心沿电场方向错开了一定距离，这就构成了电偶极子，且电偶极矩的方向与外电场方向完全相同，这种在外电场作用下产生的电偶极矩，称为\uwave{分子的感生电矩}。

由于感生电矩的存在，各个分子偶极子沿着外电场整齐链状排列，链上相邻的偶极子间正负电荷相互靠近并相互抵消，但是，链首处的正电荷，链尾处的负电荷，是无法抵消的。

这就在电介质的表面两端形成了正净电荷与负净电荷。

这种电荷就称为\uwave{极化电荷}，而电介质在外电场的作用下，电介质中的分子产生方向一致的电偶极矩的现象，称为电介质的\uwave{极化}。极化电荷的本质是极化产生的方向一致的分子偶极子的正负电荷在局部分布不均而无法抵消的净电荷。极化电荷有时也称为\uwave{束缚电荷}，这是相对于导体中可以自由移动的\uwave{自由电荷}而言的，因为在电介质中的极化电荷的实质并不是自由移动的电子，而是分子的电荷分布不均造成的，因此极化电荷是“束缚”在分子上的。\footnote{本章首的\Refx{图}{电介质对电容器的影响}中，$\sigma_f,\sigma_b$的下标就是分别代表自由电荷（\textbf{f}ree charge）和束缚电荷（\textbf{b}ound charge）。}

通过极化和极化电荷，我们就可以解释\Refx{图}{电介质对电容器的影响}中的实验现象
\begin{itemize}
    \item 由于电介质在外电场的作用下会发生极化，并在电介质的表面产生极化电荷，这就削弱了导体极板间的电场强度，从而减小了导体间的电压。但是由于电介质表面的极化电荷仅来源于表面处的分子，因此极化电荷的数量十分有限，不能完全抵消外部电场。
    \item 由于电介质中的极化电荷是束缚于分子的，是无法自由移动的，因此尽管与导体接触的电介质表面带有和导体异号的电荷，两者的电荷却并不会相互迁移发生中和。
    \begin{itemize}
        \item 电介质中的极化电荷受分子束缚，无法脱离电介质进入导体。
        \item 电容器的导体极板中的自由电荷是自由电子，也无法进入绝缘的电介质中。
    \end{itemize}
\end{itemize}
这里关于极化和极化电荷的讨论，是以匀强电场中表面与电场垂直的均匀电介质为例进行的，但这种极化原理是普遍成立的，同时我们在这里指出，只要是均匀电介质，即便其并不在匀强电场中发生均匀电极化，极化电荷也只能分布在电介质的表面，极化电荷不会分布在电介质内部，除非这是非均匀的电介质。关于“均匀电介质”和“均匀电极化”之间的差异以及结论“均匀电介质的极化电荷只能分布于电介质表面”的证明将在稍后陆续给出。

% 由此可见，电介质与导体在电场中的性质具有许多相似之处，例如对于均匀电介质与导体，两者的电荷分布都位于表面，两者的电荷分布都会在内部削弱外电场，但也有不同之处
% \begin{itemize}
%     \item 电介质通过极化产生表面的极化电荷，导体表面的电荷则来源于导体内部的自由电荷在电场作用下的定向运动。
% \end{itemize}

\begin{Figure}{无极分子的位移极化}
    \subfigure[位移极化前]
    {
        \inputfigure[scale=0.65]{Chapter09C_Figure08}
    }\quad
    % \subfigure[位移极化过程]
    % {
    %     \inputfigure[scale=0.65]{Chapter09C_Figure10}
    % }\quad
    \subfigure[位移极化后]
    {
        \inputfigure[scale=0.65]{Chapter09C_Figure09}
    }
\end{Figure}
无极分子极化的本质是电子的位移，因此无极分子的极化也称为\uwave{位移极化}。

无极分子极化的原理现在已经讨论的较为清楚了，事实上，有极分子构成的电介质在电场中也会发生类似的极化现象，但原理并不相同，下面我们就来讨论该问题。

\subsubsection{有极分子的取向极化}
对于有极分子构成的均匀电介质，如\Refx{图}{有极分子的取向极化}所示。

当电介质外没有外电场时，虽然每一分子具有固有电矩，但是由于分子无规则的热运动，在任何一块电介质中，各分子的固有电矩的空间取向概率完全相等，这就使得分子的固有电矩相互抵消，即一定区域内分子固有电矩的矢量和$\sum\vb*{p}=0$，因此在宏观上不显示出电场。

当电介质外存在外电场时，分子偶极子的正电荷和负电荷在电场的作用下受到了一对等大反向的力，这就使得分子偶极子受到了力矩的作用，而这一力矩会使得分子偶极子偏向于外电场方向，于是$\sum{\vb*{p}}\neq 0$，由于分子热运动的缘故，这种转向并不完全，因此分子偶极子最终沿外电场方向的排列不是很整齐，但总之也会在电介质两侧形成电荷。

由此可见，有极分子构成的电介质也可以发生类似的极化现象并产生极化电荷。
\begin{Figure}{有极分子的取向极化}
    \subfigure[取向极化前]
    {
        \inputfigure[scale=0.65]{Chapter09C_Figure05}
    }\quad
    \subfigure[取向极化过程]
    {
        \inputfigure[scale=0.65]{Chapter09C_Figure07}
    }\quad
    \subfigure[取向极化后]
    {
        \inputfigure[scale=0.65]{Chapter09C_Figure06}
    }
\end{Figure}
有极分子极化的本质是分子取向的一致化，因此有极分子的极化也称为\uwave{取向极化}。

有极分子与无极分子的极化原理不同，但产生的电学效果是基本一致的，两者都会产生极化电荷，两者都会使得电偶极矩的方向变得一致，因此后续讨论极化时不必再区分两者。

值得注意的是，虽然在极化时，有极分子排列不如无极分子那么整齐，有一部分电偶极矩在矢量叠加时损失了，但是，有极分子的固有电矩相较于无极分子在电场作用下产生的感生电矩要大约一个数量级，因此，有极分子的极化程度反而是比无极分子强的多的，事实上有极分子在极化时，是同时存在取向极化和位移极化的，只不过后者通常可以忽略罢了。

应当指出的是，取向极化的是通过分子的旋转实现，位移极化是通过电子的位移实现，而电子的惯性要比分子大的多，因此位移极化对电场的响应速度要比取向极化快的多，这对于静电场来说是无关紧要的，但是对于高频变化的电场，有极分子的取向极化的响应速度无法跟上电场的变化速度，此时只有惯性较小的位移极化在发挥作用，从而极化程度出现较为显著的下降，因此，通过高频电场中电介质的极化程度是否改变，就可以鉴别出电介质的类型。

\subsection{电极化强度}
通过上一小节的讨论，我们知道电介质极化的主要表现就是微观上的大量分子产生了方向相同的电偶极矩$\vb*{p}$，以至于在宏观上产生了一定的电效应，即$\sum\vb*{p}\neq 0$。

那么现在的问题就是，如何度量电介质极化程度的强弱？

如果电介质的极化程度较大，无极分子产生的感生电矩就会大一些，有极分子的固有极矩的取向也会更加一致，这就使得一定区域内分子偶极矩的矢量和$\sum\vb*{q}$更大，但是$\sum\vb*{q}$的大小受到所取区域体积的影响，因此我们要考虑单位体积的$\sum\vb*{q}$，同时为了准确的反应电介质中每一点处的极化程度，我们还需要使得取样体积$\delt V$尽可能小，即令体积$\delt V\rightarrow 0$。

基于这一想法，定义\uwave{电极化强度}
\begin{BoxDefinition}[电极化强度]
    \vb*{P}=\Lim{\delt\hspace{-1pt}V\hspace{-1pt}}{0}{\frac{\sum\vb*{p}}{\delt V}}
\end{BoxDefinition}
\Refx{定义}{电极化强度}中的$\delt{V}\rightarrow 0$是十分微妙的，一方面要求取样体积要无穷小，一方面又要求这体积内有足够多的分子，但这并不矛盾，这里的无穷小也称为\uwave{物理无穷小}，它的尺度，远小于场的宏观不均匀性，远大于场的微观不均匀性。这是合理的，因为宏观物理量的连续性本质都是微观上离散的物理量经过统计平均的结果，就像过去我们将力对时间或空间积分时，我们认为时间和空间都是连续的，不会考虑普朗克时间和普朗克长度的问题。

电极化强度反应电介质极化程度的强弱，是离散分子的电偶极矩在空间上的统计平均。

电极化强度在国际单位制中的单位是\si{C.m^{-2}}，量纲是\si{[L^{-2}TI]}。

\Refx{定义}{电极化强度}并未告诉我们电极化强度与什么因素有关，实验表明
\begin{BoxEquation}[电极化强度与场强]
    \vb*{P}=\varepsilon_0\chi_e\vb*{E}
\end{BoxEquation}
现在我们来讨论并解释\Refx{公式}{电极化强度与场强}的意义
\begin{itemize}
    \item 电极化强度$\vb*{P}$的方向与电场强度$\vb*{E}$相同，因为电场强度指向哪一方向，无极分子的感生偶极矩就指向哪一方向，有极分子的固有极矩就偏向于哪一方向。
    \item 电极化强度$\vb*{P}$的大小与电场强度$\vb*{E}$成正比，因为电场强度如果大一些，无极分子的感生极矩就相对大一些，有极分子的固有极矩的取向相对而言也就显得更为一致。
\end{itemize}
由此可见，电极化强度必然与电场强度正相关且与之方向相同。而实验进一步表明，对于通常的电介质材料，在电场强度不太大时，两者之间的正相关表现为十分良好的正比关系。

应当指出的是，该线性关系是描绘物质特性的一种尝试，然而而物质类型是非常复杂的，对于某些特殊的电介质而言，这种正比关系在相对弱的电场下就被破坏了，甚至某些电介质的电极化强度还与电场强度的变化率有关，那么在$\vb*{P},\vb*{E}$之间都不存在简单的单值对应关系
\begin{enumerate}
    \item 某些电介质在外电场撤除后仍然能保有一定的极化，这样的电介质称为\uwave{铁电体}。\footnote{铁电体的概念与磁介质中的铁磁体相对应。}
    \item 某些电介质在外电场撤除后会长期的保留极化状态。这样的电介质称为\uwave{永驻体}。\footnote{永驻体的概念与磁介质中的永磁体相对应。}
\end{enumerate}

因此\Refx{公式}{电极化强度与场强}是类似\Refx{定律}{胡克定律}的经验结论，是对于部分类型的电介质极化特性的线性近似，并不是反映电场性质，具有深刻意义的普遍性规律。

\Refx{公式}{电极化强度与场强}中的比例系数$\chi_e$称为\uwave{电极化率}，电极化率是与电介质自身的属性，表征了电介质在电场作用下极化程度的难易，电极化率较大时，无极分子更容易产生感应极矩，有极分子的固有极矩的取向更容易趋向于一致，有极分子的电极化率同时还与温度有关，因为温度较低时，分子热运动较弱，分子偶极矩的趋向也就更容易变得一致。

\Refx{公式}{电极化强度与场强}明确了“均匀电介质”和“均匀电极化”的差异
\begin{itemize}
    \item 所谓\uwave{均匀电介质}，指的是电介质处处具有相同的电极化率$\chi_e$。
    \item 所谓\uwave{均匀电极化}，指的是电介质处处具有大小和方向相同的电极化强度$\vb*{P}$，这通常是均匀电介质在匀强电场中极化时的情形（由于$\vb*{E},\chi_e$处处相同，因此$\vb*{P}$也相同）。
\end{itemize}
值得注意的是，\Refx{公式}{电极化强度与场强}中的电场强度$\vb*{E}$表示的是总场强，不仅仅包含自由电荷产生的外电场$\vb*{E}_f$，同时也包含束缚电荷产生的附加电场$\vb*{E}_b$。由于束缚电荷产生的附加电场总会削弱外电场，在实质上减小了电介质的极化，因此有时也称为\uwave{退极化场}。

\subsection{电极化强度的高斯定理}
我们知道，电极化强度是对于分子极化程度的定量描述，而分子极化会导致束缚电荷的出现，那么现在我们就要研究，电极化强度和束缚电荷间是否存在某种关系？

设想在空间中存在一块电介质$V$，现在我们在电介质中取一体积元$\dif{V}$，在这一小体积元上的电极化强度$\vb*{P}$可以视为常数。由于电极化强度本身被定义为分子偶极矩在空间上的统计平均（或者说是电偶极矩的空间密度），因此，无论体积元上大量的分子偶极子原先是如何具体分布的，无论其原先是通过位移极化还是取向极化对外显电性的，体积元$\dif{V}$产生的电效应都可以用\textbf{一个}位于体积元$\dif{V}$处且电偶极矩为$\vb*{P}\dif{V}$的电偶极子等效替代。

根据\Refx{公式}{电偶极子的电势}，体积元$\dif{V}$处的等效电偶极子产生的电势$\dif{V_b}$为
\begin{equation*}
    \dif{V_b}=\ek\frac{\vb*{P}\cdot\vu*{r}}{r^2}\dif{V}
\end{equation*}
通过三重积分，整个电介质$V$产生的电势便是
\begin{Equation}[电极化强度的高斯定理1]
    V_b=\Itnt[V]{\ek\frac{\vb*{P}\cdot\vu*{r}}{r^2}\dif{V}}
\end{Equation}
我们有必要对式\ref{式:电极化强度的高斯定理1}做一些更多的说明
\begin{Figure}{场点和源点}
    \qquad\qquad
    \inputfigure[scale=0.9]{Chapter09C_Figure11}
\end{Figure}
我们知道，三重积分的结果是一个值，但是这里计算的电势$V_b$是一个三元函数，事实上，此处是对于空间中每一个确定的场点$(x_0,y_0,z_0)$都分别计算一个三重积分，而在每一个三重积分中，场点$(x_0,y_0,z_0)$相对固定，源点$(x,y,z)$则作为变量取遍电介质$V$中的每一点。

因此式\ref{式:电极化强度的高斯定理1}中的$\vb*{r}$和其对应的单位矢量$\vu*{r}$应记为
\begin{equation*}
    \vb*{r}=(x_0-x,y_0-y,z_0-z)\qquad
    \vu*{r}=\frac{(x_0-x)\vi+(y_0-y)\vj+(z_0-z)\vk}{\sqrt{(x_0-x)^2+(y_0-y)^2+(z_0-z)^2}}
\end{equation*}
因此在每一个三重积分中，位矢$\vb*{r}$中的变量是$(x,y,z)$，因而一切对位矢$\vb*{r}$的微分运算都是对源点$(x,y,z)$的变化进行的，而此时场点$(x_0,y_0,z_0)$是相对固定的，这种场点固定源点变化的模式就是所谓的\uwave{场点坐标系}，这与我们过去所熟悉的\uwave{源点坐标系}有所不同，需要注意。

为了进一步研究式\ref{式:电极化强度的高斯定理1}，我们关注到
\begin{equation*}
    \nabla\qty(\frac{1}{r})=\nabla\frac{1}{\sqrt{(x_0-x)^2+(y_0-y)^2+(z_0-z)^2}}=
    \frac{(x_0-x)\vi+(y_0-y)\vj+(z_0-z)\vk}{\qty[(x_0-x)^2+(y_0-y)^2+(z_0-z)^2]^{\frac{3}{2}}}
\end{equation*}
即
\begin{Equation}[电极化强度的高斯定理2]
    \nabla\qty(\frac{1}{r})=\frac{\vu*{r}}{r^2}
\end{Equation}
将式\ref{式:电极化强度的高斯定理2}代入式\ref{式:电极化强度的高斯定理1}即得
\begin{Equation}[电极化强度的高斯定理3]
    V_b=\ek\Itnt[V]{\vb*{P}\cdot\nabla\qty(\frac{1}{r})\dif{V}}
\end{Equation}
根据多元微积分的理论，对于标量函数$f$与矢量函数$\vb*{A}$积的散度
\begin{equation*}
    \nabla\cdot(f\vb*{A})=\Fpif{fA_x}{x}+\Fpif{fA_y}{y}+\Fpif{fA_z}{z}
\end{equation*}
运用导数的乘积法则
\begin{equation*}
    \nabla\cdot(f\vb*{A})=f\Fpif{A_x}{x}+f\Fpif{A_y}{y}+\Fpif{A_z}{z}+A_x\Fpif{f}{x}+A_y\Fpif{f}{y}+A_z\Fpif{f}{z}
\end{equation*}
上式右侧可以重新用微分算子表示为
\begin{Equation}[散度微分恒等式]
    \nabla\cdot(f\vb*{A})=f\nabla\cdot\vb*{A}+\vb*{A}\cdot\nabla f
\end{Equation}
在式\ref{式:电极化强度的高斯定理3}中$\vb*{A}=\vb*{P}$而$f=\dfrac{1}{r}$，同时由于$\vb*{A}\cdot\nabla f=\nabla\cdot(f\vb*{A})-f\nabla\cdot\vb*{A}$，故
\begin{Equation}[电极化强度的高斯定理4]
    V_b=\ek\qty[\Itnt[V]{\nabla\cdot\qty(\frac{\vb*{P}}{r})\dif{V}}-\Itnt[V]{\frac{\nabla\cdot\vb*{P}}{r}\dif{V}}]
\end{Equation}
将式\ref{式:电极化强度的高斯定理4}的第一项通过高斯定理转换为第二类曲面积分，其中$S$为$V$的边界曲面
\begin{Equation}[电极化强度的高斯定理5]
    V_b=\ek\qty[\Isot[S]{\frac{\vb*{P}}{r}\cdot\dif{\vb*{S}}}-\Itnt[V]{\frac{\nabla\cdot\vb*{P}}{r}\dif{V}}]
\end{Equation}
将式\ref{式:电极化强度的高斯定理5}的第一项转换为第一类曲面积分
\begin{Equation}[电极化强度的高斯定理6]
    V_b=\ek\qty[\Isot[S]{\frac{\vb*{P}\cdot\vb*{n}}{r}\dif{S}}-\Itnt[V]{\frac{\nabla\cdot\vb*{P}}{r}\dif{V}}]
\end{Equation}
那么我们为何要通过这样一系列复杂的变换将电势的表达式转换为式\ref{式:电极化强度的高斯定理6}的形式？

通过\Refx{公式}{电荷元的电势}可知，如果电荷是分布在面状带电体或体状带电体上，同时，记电荷面密度为$\sigma$，记电荷体密度分别为$\rho$，那么带电体产生的电势可以分别表示为
\begin{equation*}
    V=\ek\Isot[S]{\frac{\sigma}{r}\dif{S}}\qquad
    V=\ek\Itnt[V]{\frac{\rho}{r}\dif{V}}
\end{equation*}
通过将电介质中的电偶极子产生的电势\ref{式:电极化强度的高斯定理6}与上式对比，我们可以逆向发现电介质中的束缚电荷实际上由两部分组成，一部分在电介质表面，一部分在电介质内部。

如果分别以$\sigma_b$和$\rho_b$表示这两部分电荷的电荷面密度和电荷体密度，那么
\begin{BoxPrinciple}[电极化强度与面束缚电荷]
    电极化强度与表面法向量的点积，即束缚电荷的面密度。
    \begin{BoxEq}
        \vb*{P}\cdot\vb*{n}=\sigma_b
    \end{BoxEq}
\end{BoxPrinciple}
\begin{BoxPrinciple}[电极化强度与体束缚电荷]
    电极化强度的负散度，即束缚电荷的体密度。
    \begin{BoxEq}
        \nabla\cdot\vb*{P}=-\rho_b
    \end{BoxEq}
\end{BoxPrinciple}
如果将\Refx{定律}{电极化强度与体束缚电荷}以积分形式表示\footnote{这里“闭合曲面内的束缚电荷”是指闭合曲面内的体束缚电荷，而不记闭合曲面上的面束缚电荷，事实上这两者刚好异号。}
\begin{BoxPrinciple}[电极化强度的高斯定理]
    电极化强度在闭合曲面上的负通量，即该闭合曲面内束缚电荷的总电荷量。
    \begin{BoxEq}
        \Isot[S]{\vb*{P}\cdot\dif{\vb*{S}}}=-q_b
    \end{BoxEq}
\end{BoxPrinciple}
这里可能会产生一些困惑，为何电介质在极化过程中，既会产生面束缚电荷，也会产生体束缚电荷，我们知道，束缚电荷不同于自由电荷，它并不是并不是一些带电粒子，而是分子偶极子的电荷分布不均的结果，但是这种“分布不均”实际上可以分为两种情况
\begin{enumerate}
    \item 在电介质的边界上，存在正电荷或负电荷没有可抵消的异号电荷，即面束缚电荷。
    \item 在电介质的内部中，存在正负电荷的分布密度稍有些不均的情况，即体束缚电荷。
\end{enumerate}
\begin{Figure}{面束缚电荷与体束缚电荷}
    \inputfigure[scale=0.75]{Chapter09C_Figure12}
\end{Figure}
那么对于一块电介质而言，其面束缚电荷和体束缚电荷之间存在何种关系？

根据\Refx{定律}{电极化强度与面束缚电荷}，在等式两侧同乘$\dif{S}$
\begin{equation*}
    \vb*{P}\cdot(\vb*{n}\dif{S})=\sigma_b\dif{S}
\end{equation*}
其中，左式的$\vb*{n}\dif{S}$即有向面元$\dif{\vb*{S}}$，右式$\sigma_b\dif{S}$即该面元上的电荷量$\dif{q_{b,\sigma}}$，故
\begin{equation*}
    \vb*{P}\cdot\dif{\vb*{S}}=\dif{q_{b,\sigma}}
\end{equation*}
那么对上式两侧同时求曲面积分得到
\begin{Equation}[电极化强度的通量与面束缚电荷]
    \Isot[S]{\vb*{P}\cdot\dif{\vb*{S}}}=q_{b,\sigma}
\end{Equation}
根据\Refx{定律}{电极化强度与体束缚电荷}，直接将其转化为积分形式得
\begin{Equation}[电极化强度的通量与体束缚电荷]
    \Isot[S]{\vb*{P}\cdot\dif{\vb*{S}}}=-q_{b,\rho}
\end{Equation}
关注道面束缚电荷和体束缚电荷是等量异号的，这一结果是可以预期的，因为束缚电荷的本质是分子偶极子的电荷分布不均导致的，但分子是电中性的，故分子构成的电介质整体总是电中性的\footnote{简单的说，分子偶极子总是一个正电荷带一个负电荷，两者可以在分布上不均匀，但无论如何不会有多的净电荷。}，因而束缚电荷的总和必然为零。我们将这一规律总结如下。
\begin{BoxPrinciple}[电介质的束缚电荷分布]
    电极化强度在电介质表面或电介质中任意取定的闭合曲面上的通量，

    既等于该闭合曲面上的面束缚电荷的正值，也等于该闭合曲面内的体束缚电荷的负值，

    即闭合曲面上的面束缚电荷和其内部的体束缚电荷总是等量异号，代数和为零。
\end{BoxPrinciple}
正如\Refx{定律}{电介质的束缚电荷分布}，这里的闭合曲面不一定是电介质的表面，也可以是电介质中任意取定的闭曲面，但一个常见的困惑是，我们可以理解电介质的表面有一些面束缚电荷，然而如果我们取了一个电介质内部的闭曲面，那么按照这里的理论，这个在电介质内部且是人为假象的闭曲面上也会有一些面束缚电荷分布，这看起来似乎十分荒谬。
\begin{Figure}{面束缚电荷的悖论释疑}
    \inputfigure[scale=0.75]{Chapter09C_Figure13}
\end{Figure}
在上图中，我们将电介质划分为左右两块来考虑，不难看出，在中央的假象平面上，确实出现了面束缚电荷，但这实际上并没有妨碍，这些面束缚电荷存在的原因是我们仅选取了电介质的一部分进行研究，如果我们将两块电介质合起来考虑，那么位于电介质内部的假象曲面上的面束缚电荷将会正负相互抵消，只有位于电介质边界上的面束缚电荷会被保留。

因此，虽然式\ref{式:电极化强度的通量与面束缚电荷}和式\ref{式:电极化强度的通量与体束缚电荷}的结论都是正确的，但是在\Refx{定律}{电极化强度的高斯定理}选取关于体束缚电荷的式\ref{式:电极化强度的通量与体束缚电荷}进行表述，这就是因为体束缚电荷总是在那里，不会因为我们所选取的闭曲面而变化，不会在两块电介质叠加时相互抵消。当然，这并非说面束缚电荷就都是虚幻的，那些在电介质表面的面束缚电荷与体束缚电荷一样都是极为客观真实的存在。

另外，虽然我们关于体束缚电荷讨论了这么多，但先前我们已经不加证明的指出，对于均匀电介质，即便其并不是均匀电极化，它的束缚电荷也只能分布在电介质表面，换言之\footnote{这里仍然是不加证明的先承认该结论，它的论证这里还不能完成，需要下一节的知识。}
\begin{center}
    均匀电介质，不存在体束缚电荷，只具有面束缚电荷。
\end{center}
即均匀电介质的体电荷密度处处为零，电极化强度未必是常数，但必然是无散的。

以上给出的是一种基于等效电偶极子，通过电势比较法的推证过程，是数学性较强也较为严谨的一条推导路径，下面将给出另外一种较为简单直观，但是更为常见的推导方式。
\begin{Figure}{电极化强度与束缚电荷的直观关系}
    \inputfigure[scale=1.0]{Chapter09C_Figure14}
\end{Figure}
设$\dif{\vb*{S}}$是闭合曲面上的一个有向面积元\footnote{如果所取的闭合曲面不是电介质的边界曲面，而是电介质内部的闭曲面，那么该闭曲面不应截断任何分子偶极子。}，在该面积元后沿着该处电极化强度$\vb*{P}$的方向取一长为$l$的斜柱体，其中$\vb*{l}$为该处的分子偶极矩，很显然在图示的体积为$l\dif{S}\cos\theta=\vb*{l}\cdot\dif{\vb*{S}}$灰色区域中会出现一块只存在正电荷或负电荷的薄层（这就是面束缚电荷），而这些电荷对应的分子偶极子的中心\footnote{图中以偶极矩上的红点表示。}也位于一块体积同为$\vb*{l}\cdot\dif{\vb*{S}}$的柱体区域\footnote{图中以红色虚框表示。}中，因此，如果记单位体积中的分子偶极子的数量为$n$，同时记每一分子偶极子的带电量为$q_0$，那么该区域总电荷量\footnote{这里以$\vb*{l},\dif\vb*{S}$的夹角小于$\dfrac{\pi}{2}$为例，如果$\vb*{l},\dif\vb*{S}$的夹角大于$\dfrac{\pi}{2}$会使得$\vb*{l}\cdot\dif{\vb*{S}}$的体积为负，而此时区域中的电荷刚好为负电荷。}
\begin{Equation}[电极化强度与束缚电荷的直观推导]
    \dif{q}=nq_0\vb*{l}\cdot\dif{\vb*{S}}=n\vb*{p}_0\cdot\dif\vb*{S}=\vb*{P}\cdot\dif{\vb*{S}}
\end{Equation}
其中$\vb*{p}_0$代表每一分子偶极子的电偶极矩，而显然有$\vb*{P}=n\vb*{p}_0$（单位体积的电偶极矩）。

式\ref{式:电极化强度与束缚电荷的直观推导}就说明，电极化强度在有向面积元上的通量即该面积元上的束缚电荷，将这一结论对整个闭合曲面积分，即得$\vb*{P}$在闭曲面上的通量等于其上的面束缚电荷（式\ref{式:电极化强度的通量与面束缚电荷}），由于该闭曲面内均为完整的分子，总电荷量的代数和为零，因此必然存在与面束缚电荷等量异号的体束缚电荷与之相抵消，这就得出$\vb*{P}$在闭曲面上的负通量等于其内的体束缚电荷（式\ref{式:电极化强度的通量与体束缚电荷}）。

这一推导是正确的，但是说理较为几何化，数学性有所欠缺，严谨性也有所不足。
\begin{enumerate}
    \item 该推导没有采用等效电偶极矩的观点，而是直接研究分子偶极子，而且推导基于在面积微元附近的分子偶极子的偶极矩方向均与电极化强度相同的前提，这就使得其实质上仅适用于位移极化的情况，同时这种在宏观和微观间跳跃的逻辑容易造成困惑。
    \item 该推导仅从理论上说明了闭曲面内必然存在与面束缚电荷等量异号的体束缚电荷，而不是直接导出体束缚电荷的分布，这会对理解体束缚电荷如何存在造成困难。
\end{enumerate}

\subsection{电介质电容器}
再让我们回到本章首提出的问题，先前我们已经对电介质电容器的电学效应进行了定性解释，通过\Refx{定律}{电极化强度与面束缚电荷}，现在我们可以做一些定量的研究了。
\begin{Figure}{电介质电容器}
    \inputfigure[scale=0.9]{Chapter09C_Figure15}
\end{Figure}

设电容器的正负极板充入了面密度为$\pm\sigma_f$的自由电荷，而电介质两端则在电极化的作用下产生了面密度为$\pm\sigma_b$的束缚电荷，记两者产生的场强分别为$\vb*{E}_f,\vb*{E}_b$。

设$\vb*{E}_f$的场强方向为正方向，根据\Refx{公式}{等量异号的无限大带电平面间的电场}
\begin{equation*}
    E_f=\frac{\sigma_f}{\varepsilon_0}\qquad
    E_b=-\frac{\sigma_b}{\varepsilon_0}
\end{equation*}
设$\vb*{E}$为电介质电容器中的总场强，则
\begin{equation*}
    E=E_f+E_b=E_f-\frac{\sigma_b}{\varepsilon_0}
\end{equation*}
这里我们的目的是要研究电容器置入电介质前后的电压关系，这相当于就是要研究电容器置入电介质前后的电场强度的关系，即$E,E_f$间的关系，但$E$的表达式还包含未知量$\sigma_b$。

在上式中代入\Refx{定律}{电极化强度与面束缚电荷}
\begin{equation*}
    E=E_f-\frac{\vb*{P}\cdot\vb*{n}}{\varepsilon_0}=E_f-\frac{P}{\varepsilon_0}
\end{equation*}
在上式中继续代入\Refx{公式}{电极化强度与场强}
\begin{equation*}
    E=E_f-\frac{\varepsilon_0\chi_eE}{\varepsilon_0}
\end{equation*}
即有
\begin{equation*}
    E=E_f-\chi_eE
\end{equation*}
移项解出总场强$E$，即得
\begin{BoxEquation}[电介质电容器场强的电极化率表达]
    E=\frac{E_f}{1+\chi_e}
\end{BoxEquation}
这就表明电场强度在电介质中会被削弱$1+\chi_e$倍，若将这个倍数记为$\varepsilon_r$，那么
\begin{BoxEquation}[电介质电容器场强的电容率表达]
    E=\frac{E_f}{\varepsilon_r}
\end{BoxEquation}
这里的$\varepsilon_r$称为电介质的\uwave{相对电容率}或\uwave{相对介电常数}，我们稍后会看到这个名称的意义
\begin{BoxDefinition}[相对电容率]
    \varepsilon_r=1+\chi_e
\end{BoxDefinition}
相对电容率和电极化率在国际单位制中均为无量纲量。

根据\Refx{公式}{电介质电容器场强的电容率表达}，在置入电介质后电场强度被削弱了$\varepsilon_r$倍，由于匀强电场中电压$U=Ed$，因此，在置入电介质后电压也被削弱了$\varepsilon_r$倍，即
\begin{Equation}[电介质电容器的电压]
    \frac{U}{U_f}=\frac{E\cdot d}{E_f\cdot d}=\frac{1}{\varepsilon_r}
\end{Equation}
根据\Refx{定义}{电容}，在置入电介质后电容将增大至原先的$\varepsilon_r$倍，即
\begin{Equation}[电介质电容器的电容]
    \frac{C}{C_f}=\frac{U^{-1}\cdot q}{U_f^{-1}\cdot q}=\varepsilon_r
\end{Equation}
真空中并没有可极化的分子，因此，真空的电极化率$\chi_e=0$，真空的相对电容率$\varepsilon_r=1$，而相比之下，电介质的电极化率$\chi_e>0$，电介质的相对电容率$\varepsilon_r>1$，结合式\ref{式:电介质电容器的电容}可知电介质电容器的电容是对应真空电容器的$\varepsilon_r$倍，这就是相对电容率的含义，\textbf{即相对于真空电容器的电容比率}，由于这一性质与电介质有关，因此相对电容率也称为相对介电常数。

值得注意的是，尽管相对电容率$\varepsilon_r$的概念从电介质电容器中产生，但它在实质上只是一个简单的定义$\varepsilon_r=1+\chi_e$，我们之后也会在很多“不太电容”的地方用到$\varepsilon_r$的表示。

在上一节中我们曾提到过，电容器的极板面积一定时，电容器的电容与耐压是矛盾的，通过减小极板间距增大电容必然会导致耐压下降。但是这里我们看到，通过在极板间填充电介质也可以在不改变耐压的前提下增大电容，如果选取了恰当的电介质，电容可以增大数十倍到数百倍，性能得到显著提升，现在的电子线路中的电容器基本均为电介质电容器。